% GENERATED FILE. Edit canonical lessons, then run npm run generate:books.
% canonical-source-hash: fnv1a64:85f39dbf8b16b03f
% canonical-lessons: GE-C156-rezept, GE-C156-schmerz, GE-C156-erkaeltung, GE-C156-husten, GE-C156-grippe

\chapter{A prescription, A pain, A cold, A cough, The flu}
\label{ch:ge-a-prescription-a-pain-a-cold-a-cough-the-flu}
\hlchaptermodality{156}

\begin{chapteropening}
\textbf{By the end of this chapter:} \emph{I can say a prescription, a pain, a cold, a cough and the flu.}

Complete the last lesson of chapter 156: I can say a prescription, a pain, a cold, a cough and the flu.
\end{chapteropening}

\section[das Rezept]{das Rezept — a prescription, a recipe}

\label{lesson:GE-C156-rezept}



\begin{quote}
Before the new one: say the German for a traffic jam, then the German for an accident.
\end{quote}

\subsection*{You'll want to know: das Rezept}
\textbf{das Rezept} — \textquotedblleft{}a prescription, a recipe\textquotedblright{}.

\begin{grammarlens}[title={What you've built}]
One more word for this chapter.
\end{grammarlens}

\subsection*{Your turn}
\begin{itemize}
\raggedright
  \item \emph{Say it:} \emph{das Rezept}
  \item \emph{Say it:} \emph{das Rezept}, once more
  \item \emph{Say it:} say \emph{das Rezept}
  \item \emph{Recall:} say the German for a traffic jam, then the German for an accident, then say \emph{das Rezept} again
\end{itemize}

\begin{tcolorbox}[breakable,colback=teal!4,colframe=teal!35!black,title={Before you move on}]
What does das Rezept mean? (\textquotedblleft{}A prescription, a recipe\textquotedblright{}.) Say it once more.
\end{tcolorbox}
\section[der Schmerz]{der Schmerz — a pain}

\label{lesson:GE-C156-schmerz}



\begin{quote}
Before the new one: say the German for an accident, then the German for a prescription.
\end{quote}

\subsection*{You'll want to know: der Schmerz}
\textbf{der Schmerz} — \textquotedblleft{}a pain\textquotedblright{}.

\begin{grammarlens}[title={What you've built}]
One more word for this chapter.
\end{grammarlens}

\subsection*{Your turn}
\begin{itemize}
\raggedright
  \item \emph{Say it:} \emph{der Schmerz}
  \item \emph{Say it:} \emph{der Schmerz}, once more
  \item \emph{Say it:} say \emph{der Schmerz}
  \item \emph{Recall:} say the German for an accident, then the German for a prescription, then say \emph{der Schmerz} again
\end{itemize}

\begin{tcolorbox}[breakable,colback=teal!4,colframe=teal!35!black,title={Before you move on}]
What does der Schmerz mean? (\textquotedblleft{}A pain\textquotedblright{}.) Say it once more.
\end{tcolorbox}
\section[die Erkältung]{die Erkältung — a cold}

\label{lesson:GE-C156-erkaeltung}



\begin{quote}
Before the new one: say the German for a prescription, then the German for a pain.
\end{quote}

\subsection*{You'll want to know: die Erkältung}
\textbf{die Erkältung} — \textquotedblleft{}a cold\textquotedblright{}.

\begin{grammarlens}[title={What you've built}]
One more word for this chapter.
\end{grammarlens}

\subsection*{Your turn}
\begin{itemize}
\raggedright
  \item \emph{Say it:} \emph{die Erkältung}
  \item \emph{Say it:} \emph{die Erkältung}, once more
  \item \emph{Say it:} say \emph{die Erkältung}
  \item \emph{Recall:} say the German for a prescription, then the German for a pain, then say \emph{die Erkältung} again
\end{itemize}

\begin{tcolorbox}[breakable,colback=teal!4,colframe=teal!35!black,title={Before you move on}]
What does die Erkältung mean? (\textquotedblleft{}A cold\textquotedblright{}.) Say it once more.
\end{tcolorbox}
\section[der Husten]{der Husten — a cough}

\label{lesson:GE-C156-husten}



\begin{quote}
Before the new one: say the German for a pain, then the German for a cold.
\end{quote}

\subsection*{You'll want to know: der Husten}
\textbf{der Husten} — \textquotedblleft{}a cough\textquotedblright{}.

\begin{grammarlens}[title={What you've built}]
One more word for this chapter.
\end{grammarlens}

\subsection*{Your turn}
\begin{itemize}
\raggedright
  \item \emph{Say it:} \emph{der Husten}
  \item \emph{Say it:} \emph{der Husten}, once more
  \item \emph{Say it:} say \emph{der Husten}
  \item \emph{Recall:} say the German for a pain, then the German for a cold, then say \emph{der Husten} again
\end{itemize}

\begin{tcolorbox}[breakable,colback=teal!4,colframe=teal!35!black,title={Before you move on}]
What does der Husten mean? (\textquotedblleft{}A cough\textquotedblright{}.) Say it once more.
\end{tcolorbox}
\section[die Grippe]{die Grippe — the flu}

\label{lesson:GE-C156-grippe}



\begin{quote}
Before the new one: say the German for a cold, then the German for a cough.
\end{quote}

\subsection*{You'll want to know: die Grippe}
\textbf{die Grippe} — \textquotedblleft{}the flu\textquotedblright{}.

\begin{grammarlens}[title={What you've built}]
One more word for this chapter.
\end{grammarlens}

\subsection*{Your turn}
\begin{itemize}
\raggedright
  \item \emph{Say it:} \emph{die Grippe}
  \item \emph{Say it:} \emph{die Grippe}, once more
  \item \emph{Say it:} say \emph{die Grippe}
  \item \emph{Recall:} say the German for a cold, then the German for a cough, then say \emph{die Grippe} again
\end{itemize}

\begin{tcolorbox}[breakable,colback=teal!4,colframe=teal!35!black,title={Before you move on}]
What does die Grippe mean? (\textquotedblleft{}The flu\textquotedblright{}.) Say it once more.
\end{tcolorbox}
